\documentclass[UTF8,fontset=none,11pt,oneside,openany]{ctexbook}

\usepackage{iftex}
\usepackage{fontspec}
\ifXeTeX
  \usepackage{xeCJK}
\fi
\defaultfontfeatures{Ligatures=TeX}
\setmainfont{texgyretermes-regular.otf}[
  BoldFont=texgyretermes-bold.otf,
  ItalicFont=texgyretermes-italic.otf,
  BoldItalicFont=texgyretermes-bolditalic.otf
]
\setsansfont{texgyreheros-regular.otf}
\setmonofont{lmmono10-regular.otf}[ItalicFont=lmmono10-italic.otf]
\setCJKmainfont{FandolSong-Regular.otf}[
  BoldFont=FandolSong-Bold.otf,
  ItalicFont=FandolKai-Regular.otf,
  BoldItalicFont=FandolSong-Bold.otf
]
\setCJKsansfont{FandolHei-Regular.otf}
\setCJKmonofont{FandolFang-Regular.otf}
\ifXeTeX
  \xeCJKsetup{PunctStyle=kaiming,CheckSingle=true}
  \XeTeXlinebreaklocale "zh"
  \XeTeXlinebreakskip=0pt plus 1pt
\fi

\usepackage{amsmath}
\ifLuaTeX
  \usepackage{unicode-math}
  \setmathfont{texgyretermes-math.otf}
  \providecommand{\square}{\mdlgwhtsquare}
\else
  \usepackage{amssymb,mathrsfs}
\fi
\ifLuaTeX
  \renewcommand{\leadsto}{\mathrel{\rightsquigarrow}}
\fi

\def\sh{\mathcal}
\def\bb{\mathbf}
\def\cat{\mathtt}
\def\leq{\leqslant}
\def\geq{\geqslant}
\def\setmin{-}
\def\rad{\mathfrak{r}}
\def\nilrad{\mathfrak{N}}
\def\emp{\varnothing}
\def\vphi{\phi}
\def\HH{\mathrm{H}}
\def\CHH{\check{\HH}}
\def\RR{\mathrm{R}}
\def\LL{\mathrm{L}}
\def\dual#1{{#1}^\vee}
\def\kres{k}
\def\C{\cat{C}}
\def\op{^\cat{op}}
\def\Set{\cat{Set}}
\def\CHom{\cat{Hom}}
\def\supertilde{{\,\widetilde{\,}\,}}
\def\GL{\bb{GL}}
\def\bbGamma{\boldsymbol\Gamma}
\def\red{\mathrm{red}}
\def\rg{\operatorname{rg}}
\def\gr{\operatorname{gr}}
\def\Gr{\operatorname{Gr}}
\def\Hom{\operatorname{Hom}}
\def\Proj{\operatorname{Proj}}
\def\Tor{\operatorname{Tor}}
\def\Ext{\operatorname{Ext}}
\def\Supp{\operatorname{Supp}}
\def\Ker{\operatorname{Ker}}
\def\Im{\operatorname{Im}}
\def\Coker{\operatorname{Coker}}
\def\Spec{\operatorname{Spec}}
\def\Spm{\operatorname{Spm}}
\def\Spf{\operatorname{Spf}}
\def\grad{\operatorname{grad}}
\def\dimc{\operatorname{dimc}}
\def\codim{\operatorname{codim}}
\def\id{\operatorname{id}}
\def\Der{\operatorname{Der}}
\def\Diff{\operatorname{Diff}}
\def\Hyp{\operatorname{Hyp}}
\def\Tr{\operatorname{Tr}}

\usepackage{graphicx}
\usepackage{tikz-cd}
\ifLuaTeX
  \usepackage{luatex85}
  \input{glyphtounicode.tex}
  \pdfgentounicode=1
\fi
\usepackage[all]{xy}
\ifLuaTeX
  \makeatletter
  \AtBeginDocument{%
    \edef\xP@preclw{0.43799pt}%
    \edef\xP@lw{\xP@dim\xP@preclw}%
  }
  \makeatother
\fi
\usepackage[shortlabels]{enumitem}
\usepackage{marginnote}
\newcommand{\SourceCorrectionNote}[1]{\footnote{#1}}
\ifLuaTeX\else\usepackage{microtype}\fi
\usepackage{fancyhdr}
\usepackage[
  a4paper,
  inner=26mm,
  outer=24mm,
  top=24mm,
  bottom=24mm,
  headheight=14pt,
  headsep=8mm,
  footskip=12mm
]{geometry}
\usepackage[
  unicode=true,
  colorlinks=true,
  linkcolor=black,
  citecolor=black,
  urlcolor=blue,
  hypertexnames=false
]{hyperref}

\hypersetup{
  pdftitle={代数几何原理 IV-1：概形与概形态射的局部研究（大陆简体中文工作检查点）},
  pdfauthor={Alexander Grothendieck；Jean Dieudonné（协作编写）},
  pdfsubject={依据 NUMDAM 法文印本与法文外交式 TeX 制作的大陆简体中文版本},
  pdfkeywords={代数几何；概形；态射；EGA；简体中文}
}

\AtBeginDocument{\fontsize{11pt}{16.5pt}\selectfont}
\setlength{\parindent}{2em}
\setlength{\parskip}{0pt}
\setlength{\emergencystretch}{2em}
\flushbottom

\pagestyle{fancy}
\fancyhf{}
\fancyhead[L]{\small 代数几何原理 IV-1}
\fancyhead[R]{\small 概形与态射的局部研究}
\fancyfoot[C]{\thepage}
\renewcommand{\headrulewidth}{0.4pt}
\renewcommand{\footrulewidth}{0pt}
\fancypagestyle{plain}{%
  \fancyhf{}%
  \fancyhead[L]{\small 代数几何原理 IV-1}%
  \fancyhead[R]{\small 概形与态射的局部研究}%
  \fancyfoot[C]{\thepage}%
  \renewcommand{\headrulewidth}{0.4pt}%
  \renewcommand{\footrulewidth}{0pt}%
}

\renewcommand{\contentsname}{目录}
\renewcommand{\bibname}{参考文献}
\renewcommand{\figurename}{图}
\renewcommand{\tablename}{表}
\renewcommand{\thesection}{\arabic{section}}
\renewcommand{\thesubsection}{\thesection.\arabic{subsection}}

\newenvironment{env}[1][]{%
  \par\medskip\noindent\textbf{(#1)}\ %
}{\par}
\newenvironment{definition}[1][]{%
  \par\medskip\noindent\textbf{定义 (#1).}\ %
}{\par}
\newenvironment{lemma}[1][]{%
  \par\medskip\noindent\textbf{引理 (#1).}\ %
}{\par}
\newenvironment{theorem}[1][]{%
  \par\medskip\noindent\textbf{定理 (#1).}\ %
}{\par}
\newenvironment{proposition}[1][]{%
  \par\medskip\noindent\textbf{命题 (#1).}\ %
}{\par}
\newenvironment{corollary}[1][]{%
  \par\medskip\noindent\textbf{推论 (#1).}\ %
}{\par}
\newenvironment{remark}[1][]{%
  \par\medskip\noindent\textbf{注 (#1).}\ %
}{\par}
\newenvironment{remarks}[1][]{%
  \par\medskip\noindent\textbf{注 (#1).}\ %
}{\par}
\newenvironment{notation}[1][]{%
  \par\medskip\noindent\textbf{记号 (#1).}\ %
}{\par}
\newenvironment{example}[1][]{%
  \par\medskip\noindent\textbf{例 (#1).}\ %
}{\par}
\newenvironment{examples}[1][]{%
  \par\medskip\noindent\textbf{例 (#1).}\ %
}{\par}
\newenvironment{proof}{%
  \par\smallskip\noindent\textit{证明.}\ %
}{\hfill$\square$\par}

\makeatletter
\newcommand{\oldpage}[2][]{%
  \ifcsname ega@sourcepage@IVone@#2\endcsname
  \else
    \expandafter\gdef\csname ega@sourcepage@IVone@#2\endcsname{1}%
    \hypertarget{ega-source-page-IV1-#2}{}%
  \fi
  \marginnote{\scriptsize\textbf{#1}~|~#2}\ignorespaces}
\makeatother

\let\EGAOriginalLabel\label
\renewcommand{\label}[1]{\phantomsection\EGAOriginalLabel{#1}}
